\section{Elastic Beanstalk}

\subsection{AWS Elastic Beanstalk Overview}
\begin{itemize}
    \item Elastic Beanstalk is a developer centric view of deploying an application on AWS
    \item It uses all the components we've seen before: EC2, Auto Scaling Group, Elastic Load Balancers, RDS, etc.......
    \item We still have full control over the configuration of each component
    \item Beanstalk is free but you pay for the underlying instances
\end{itemize}

\subsection{Elastic Beanstalk}

\begin{itemize}
    \item Support for many platforms:
    \begin{itemize}
        \item Go
        \item Java SE
        \item Java with Tomcat
        \item .NET on Window Server with IIS
        \item Node.js
        \item PHP
        \item Python
        \item Ruby
        \item Packer Builder
        \item Single Container Docker
        \item Multi-container Docker
        \item Preconfigured Docker
    \end{itemize}
\end{itemize}

- Beanstalk is great to "platform" your application from on-premises to the cloud

\begin{itemize}
    \item Managed service
    \begin{itemize}
        \item Instance configuration / OS is handled by Beanstalk
        \item Deployment strategy is configurable by performed by Elastic Beanstalk
    \end{itemize}
    \item Just the application code is the responsibility of the developer
    \item Three architecture models:
    \begin{itemize}
        \item Single instance deployment: good for developers
        \item LB + ASG: great for production or pre-production web applications
        \item ASG only: great for non-web apps in production (workers. etc......)
    \end{itemize}
\end{itemize}

\subsection{Beanstalk Environments}
% TODO add diagrams

\subsection{Web Server vs Worker Environment}
\begin{itemize}
    \item If your application performs tasks that are long to complete, offload these tasks to a dedicated worker environment
    \item Decoupling you application into two tiers is common
    \item Example: processing a video, generating a zip file, etc
    \item You can define periodic task in a file cron.yaml
\end{itemize}

%TODO add diagrams

\subsection{Elastic Beanstalk Deployment Blue / Green}
%TODO add diagrams
\begin{itemize}
    \item Not a "direct feature" of Elastic Beanstalk
    \item Zero downtime and release facility
    \item Create a new "stage" environment and deploy v2 there
    \item The new environment (green) can be validated independently and roll back if there are issues
    \item Route 53 can be setup using weighted policies to redirect a little bit of traffic to the stage environment
    \item Using Beanstalk "swap URLs" (DNS swap) when done with the environment test
\end{itemize}


\section{CodeDeploy}

\subsection{AWS CodeDeploy}
%TODO diagram
\begin{itemize}
    \item We want to deploy out application automatically to many EC2 instances
    \item These instances are no managed by Elastic Beanstalk
    \item There are several ways to handle deployment using open source tools (Ansible, Terraform, Chef, Puppet, etc......)
    \item We can use the managed Service AWS CodeDeploy
    \item CodeDeploy can deploy to: EC2, ASG, ECS and Lamabda
\end{itemize}

\subsection{CodeDeploy to EC2}
%TODO add diagram
\begin{itemize}
    \item Define how to deploy the application using appspec.yml + deployment strategy
    \item Will do in-place update your fleet of EC2 instances
    \item Can use hooks to verify the deployment after each deployment phase
\end{itemize}

\subsection{CodeDeploy to ASG}
%TODO add diagram
\begin{itemize}
    \item In place updates:
    \begin{itemize}
        \item Updated current existing EC2 instances
        \item Instances newly created by an ASG will also get automated deployments
    \end{itemize}
    \item Blue / green deployment:
    \begin{itemize}
        \item A new auto-scaling group is created (settings are copied)
        \item Choose how long to keep the old instances
        \item Must be using an ELB
    \end{itemize}
\end{itemize}

\subsection{CodeDeploy to AWS Lambda}
%TODO add diagram
\begin{itemize}
    \item Traffic Shifting feature
    \item Pre and Post traffic hooks features to validate deployment (before the traffic shift starts and after it ends)
    \item Easy and automated rollback using CloudWatch Alarms
    \item SAM framework natively uses CodeDeploy
\end{itemize}

\subsection{CodeDeploy to ECS}
% TODO add diagram
\begin{itemize}
    \item Support for Blue/Green deployments for Amazon ECS and AWS Fargate
    \item Setup is done within the ECS service
    \item A new task set is created and traffic is re-routed to the new task test
    \item Then if everything is stable for X minutes, the old task set is terminated (so you have time to notice the issues)
    \item Supports Canary deployments (Canary10Percent5Minute)
\end{itemize}


\section{CloudFormation}
\begin{itemize}
    \item Infrastructure as code (IaC) in AWS
    \item Portability of stacks across multiple accounts and regions
    \item Backbone of the elastic beanstalk service
    \item Backbone of the Service Catalog service
    \item Backbone of the SAM (Serverless Application Model) framework
    \item Must-know service as a developer / sysops / devops
\end{itemize}

\subsection{Retaining Data on Deletes}
\begin{itemize}
    \item You can put a DeletionPolicy on any resource to control what happens when the CloudFormation template is deleted
    \item DeletionPolicy=Retain
    \begin{itemize}
        \item Specify on resources to preserve / backup in case of CloudFormation deleted
        \item To keep a resource, specify Retain (works for any resource / nested stack)
    \end{itemize}
    \item DeletionPolicy=Snapshot:
    \begin{itemize}
        \item EBS Volume, Elasticache Cluster, ElastiCache ReplicationGroup
        \item RDS DBInstance, RDS DBCluster, Redshift Cluster
    \end{itemize}
    \item DeletePolicy=Delete (default behaviour)
    \begin{itemize}
        \item Note: for AWS::RDS::DBCluster resources, the default policy is Snapshot
        \item Note: to delete an S3 bucket, you first empty the bucket of its content
    \end{itemize}
\end{itemize}

\subsection{CloudFormation Customer Resources (Lambda)}
%TODO diagram
\begin{itemize}
    \item You can define a Custom Resource in CloudFormation to address any of these use cases
    \item An AWS resource is not yet supported (new service for example)
    \item An on-premises resource
    \item Emptying an S3 bucket before being deleted
    \item Fetch an AMI id
\end{itemize}

\subsection{CloudFormation - StackSets}
%TODO diagram add
\begin{itemize}
    \item Create, update or delete stacks across multiple accounts and regions with a single operation
    \item Administrator account to create StackSets
    \item Trusted accounts to create, update, delete stack instances from StackSets
    \item When you updated a stack set, all associated stack instances are updated through all accounts and regions
    \item Enabled Automatic Deployment feature to autmatically deploy to accounts in AWS Organisations or OUs
\end{itemize}

\subsection{CloudFormation Drift Overview}
%TODO add diagram
\begin{itemize}
    \item CloudFormation allows you to create infrastructure
    \item But it doesn't protect you against manual configuration changes
    \item How do we know if our resources have drifted? -> Use cloud formation drift
    \item Detect drift on an entire stack or on individual resource within a stack
\end{itemize}

\subsection{CloudFormation - Integration with Secrets Manager}
%TODO see screenshot

\subsection{CloudFormation - Resource Import}
% TODO add diagram
\begin{itemize}
    \item Import existing resources into existing / new stack s
    \item You don't need to delete and re-create the resources as part of a CloudFormation stack
    \item During import operation, you'll ned
    \item A template that describes the entire stack (original stack resources and target resources to import)
    \item A Unique identifier for each target resource (ex BucketName for S3 buckets)
    \item Each resource to import must have a DeletionPolicy attribute (any value) and Identifier
    \item Can't import the same resource into multiple stacks
\end{itemize}


\section{Service Catalog}

\subsection{AWS Service Catalog}
\begin{itemize}
    \item Users that are new to AWS have too many options and may create stacks that are not compliant / in line with the rest of the organisation
    \item Some users just want a quick self-service portal to launch a set of authorised products pre-defined by admins
    \item Includes: virtual machines, databases, storage options, etc.........
    ---> AWS Service Catalog use case
\end{itemize}

\subsection{Service Catalog Diagram}
%TODO add diagram

\subsection{AWS Service Catalog}
\begin{itemize}
    \item Create and manage catalogs of IT services that are approved on AWS
    \item The "products" are CloudFormation templates
    \item Ex:Virtual machine images, Servers, Software, Databases, Regions, IP address ranges
    \item CloudFormation helps ensure consistency and standardisation by Admins
    \item They are assigned to Portfolios (teams)
    \item Team are presented a self-service portal where they can launch the products
    \item All the deployed products are centrally managed deployed services
    \item Help with governance, compliance and consistency
    \item Can give user access to launching products withouth requiring deep AWS knowledge
    \item Integrations with "self-service portals" such as ServiceNow
\end{itemize}


\section{SAM - Serverless Application Model}
\begin{itemize}
    \item SAM = Serverless Application Model
    \item Framework for developing and deploying serverless applications
    \item All the configuration is YAML code. Examples:
    \item Lambda Functions (AWS::Serverless::Function)
    \item DynamoDB tables (AWS::ServerLess::SimpleTable)
    \item API Gateway (AWS::Serverless::API)
    \item StepFunction - State Machine (AWS::Serverless::StateMachine)
    \item SAM can help you to run Lambda, API Gateway, DynamoDB locally
    \item SAM can use CodeDeploy to deploy Lambda functions (traffic shifting)
    \item Leverages CloudFormation in the backend
\end{itemize}

\subsection{CICD Architecture for SAM}
%TODO add diagram


\section{AWS CDK - Cloud Development Kit}
%TODO add diagram
\begin{itemize}
    \item Define your cloud infrastructure using a familiar language:
    \item JavaScript/ TypeScript, Python, Java and .NET
    \item The code is "compiled" into a CloudFormation template (JSON/YAML)
    \item You can therefore deploy infrastructure and application runtime code together
    \item Great for Lambda functions
    \item Great for Docker containers in ECS / EKS
\end{itemize}

\subsection{CDK Example}
%TODO see example


\section{AWS Systems Manger - SSM}

\subsection{AWS Systems Manager Overview}
\begin{itemize}
    \item Helps you manage your EC2 and on-premises system at scal
    \item Get operational insights about the state of your infrastructure
    \item Easily detect problems
    \item Patching automation for enhancd compliance
    \item Works for both Windows and Linux OS
    \item Integrated with CloudWatch metrics / dashboards
    \item Integrated with AWS Config
    \item Free service
\end{itemize}

\subsection{How Systems Manger Works}
%TODO see diagram
\begin{itemize}
    \item We need to install the SSM agent onto the systems we controll
    \item Installed by default on Amazon Linux AMI and some Ubuntu AMI
    \item If an instance can't be controlled with System Manger, it's probably an issue with the SSM agent!
    \item Make sure the EC2 instances have a proper IAM role to allow System Manager actions
\end{itemize}

\subsection{AWS Systems Manager Run Command}
%TODO add diagram
\begin{itemize}
    \item Execute a document (script) or just run a command
    \item Run command across multiple instances (using resources groups)
    \item Rate Control / Error Control
    \item Integrated with IAM and CloudTrail
    \item No need for SSH
    \item Results in the console
\end{itemize}

\subsection{AWS Systems Manager - Send Command before an ASG Instance is Terminated}
%TODO add diagrma
\begin{itemize}
    \item Perform any action before the ASG terminates an EC2 instance
    \item Create a ASG lifecycle hook that puts the instance in Terminating:Wait state
    \item Monitor the Terminating:Wait state using EventBridge
    \item Trigger a SSM automation document to perform the actions on the instance before termination
\end{itemize}

\subsection{Systems Manager Patch Managers - Steps}
%TODO add diagram
\begin{enumerate}
    \item Define a patch baseline to use (or multiple if you have multiple environments)
    \item Define patch groups: define based on tags, example different environments (dev, test, prod) --- use tag Patch Group
    \item Define Maintenance Windows (schedule, duration, registered targets/ patch groups and registered tasks)
    \item Add the AWS-RunPatchBaseline Run Command as part of the registered tasks of the Maintenance Window (works cross platform Linux and Windows)
    \item Define Rate control (concurrency and error threshold) for the task
    \item Monitor Patch Compliance using SSM inventory
\end{enumerate}

\subsection{Systems Manager / Session Manager}
%TODO add diagram
\begin{itemize}
    \item Allows you to start a secure shell on your EC2 and on-premises servers
    \item Access through AWS Console, AWS ClI or Session Manager SDK
    \item Does not need SSH access, bastion hosts or SSH keys
    \item Supports Linux, macOS and Windows
    \item Log connections to your instances and executed commands
    \item Session log data can be sent to S3 or CloudWatch Logs
    \item CloudTrail can intercept StartSession events
\end{itemize}

\subsection{Systems Manager OpsCenter}
\begin{itemize}
    \item Resolve Operational Issues (OpsItems) related to AWS resources
    \item OpsItems - issues, events and alerts
    \item Aggregates information to resolve issue on each OpsItems such as:
    \item AWS Config changes and relationships
    \item CloudTrail Logs
    \item CloudWatch Alarms
    \item CloudFormation Stack Information
    \item Provides Automation Runbook that you can resolve issues
    \item EventBridge or CloudWatch Alarms can create OpsItems
\end{itemize}


\section{AWS Cloud Map}

\subsection{AWS Cloud Map}
%TODO add diagram
\begin{itemize}
    \item A fully managed resource discovery service
    \item Creates a map of the backend services / resources that your application depend on
    \item You register your application components, their locations, attributes and health status with AWS Cloud Map
    \item Integrated health checking (stop sending traffic to unhealthy endpoints)
    \item Your applications can query AWS Cloud Map using AWS SDK, API or DNS
\end{itemize}
